\clearpage\SourceLeaf{134}%
\pdfbookmark[1]{Über eine Aufgabe der sphärischen Astronomie}{sphaerisch1808}
\fancyhead[L]{\small ÜBER EINE AUFGABE DER SPHÄRISCHEN ASTRONOMIE}
\begin{center}
{\small ÜBER EINE AUFGABE DER SPHÄRISCHEN ASTRONOMIE.\par}
\medskip\hrule\medskip
{\small Monatliche Correspondenz zur Beförderung der Erd- und Himmels-Kunde,\par
herausgegeben vom Freiherrn \textsc{von Zach}. Oktober 1808.\par}
\medskip\hrule
\end{center}\medskip
Die tägliche Bewegung der Himmelskörper wirft eine große Vielfalt von Problemen auf, bei denen es um die Beziehungen zwischen Stundenwinkeln, Höhen, Azimuten, den Positionen der Himmelskörper und der Polhöhe geht. \textsc{Maupertuis} und andere Astronomen haben sich mit mehreren dieser Probleme beschäftigt. Sie sind allerdings von sehr unterschiedlichem und zum Teil sehr geringem praktischem Nutzen. Solche Aufgaben hat man vor allem Seeleuten und Reisenden sowie Beobachtern empfohlen, die nur über wenige und weniger vollkommene Instrumente verfügen. Sie sollen bei den beiden wichtigsten Beobachtungen helfen, der Zeit- und der Ortsbestimmung. Für Astronomen, denen die besten Uhren, Meridianteleskope, Mauerquadranten, Vollkreise und Zenitsektoren zur Verfügung stehen, sind solche Hilfsmethoden entbehrlich. Sie können die Zeit zu jeder Stunde mit größter Genauigkeit und Bequemlichkeit bestimmen. Welche Methoden sich am besten eignen, die geografische Lage ihres Beobachtungsorts festzulegen, ist für sie keine offene Frage mehr. Natürlich bleibt auch die Bestimmung der Sternpositionen und der Positionen der Sonne allein diesen Astronomen vorbehalten. Ein Beobachter mit schlechteren Instrumenten wird diese Positionen immer von ihnen übernehmen, wenn er sie braucht, und sie nicht aus seinen eigenen, weniger genauen Beobachtungen ableiten wollen. Daher ist es zum Beispiel praktisch völlig wertlos, aus drei beobachteten Sternhöhen
\SourceLeaf{135}%
zugleich die Polhöhe, die Deklination und die Kulminationszeit des Sterns zu bestimmen. Die einzige, kaum vorstellbare Ausnahme wäre, dass man nur einen einzigen Stern beobachten könnte, dessen Position noch nicht gut bestimmt ist. Jenes Verfahren könnte nur dann brauchbare Ergebnisse liefern, wenn die Beobachtungen zeitlich sehr weit auseinanderlägen und zugleich sowohl der Gang der Uhr währenddessen als auch die gemessenen Höhen selbst sehr genau wären. Selbst dann lassen sich Polhöhe und Zeit immer \textit{viel} genauer aus zwei dieser Beobachtungen ableiten, wenn man die Deklination des Sterns als gegeben ansieht.

Eine der nützlichsten Aufgaben überhaupt für Seeleute und reisende Beobachter ist es, aus zwei beobachteten Höhen zweier Sterne und den zugehörigen Uhrzeiten den Uhrstand und die Polhöhe zu bestimmen. Dabei gelten die Rektaszensionen und Deklinationen der Sterne als gegeben, und die Uhr geht entweder nach Sternzeit oder ihr Gang während der Beobachtungen wird als bekannt vorausgesetzt. Die Sterne und Höhen sind frei wählbar. Die einzige Bedingung ist, dass die Vertikalkreise, in denen die Höhen gemessen werden, im Zenit einen Winkel bilden, der weder zu spitz ist noch zu nahe bei $180^0$ liegt. So kann man die beiden Höhen leicht innerhalb weniger Minuten messen. Das ist zweifellos von größter Bedeutung für den Seefahrer, dessen Standort sich auf dem Meer ständig rasch ändert, ebenso wie für den reisenden Beobachter an Land, dessen Zeit begrenzt ist, der nur wenige Minuten lang günstiges Wetter hat oder der sich auf den Gang seiner Uhr nicht lange verlassen kann. Die naheliegende direkte Lösung dieser Aufgabe beruht auf der Berechnung von drei sphärischen Dreiecken und ist freilich etwas umständlich. Die indirekte Methode lässt sich sehr bequem und flexibel gestalten. Ich behalte mir vor, ein andermal darauf zurückzukommen.

Wenn man Sternhöhen mit einem Reflexionsinstrument und einem künstlichen Horizont beobachtet, ist es anfangs etwas schwierig, die beiden Bilder ins Gesichtsfeld zu bringen. Man erwirbt darin aber bald Übung, zumal wenn man sich auf die helleren Sterne beschränkt. Eine ganz grobe Vorausberechnung der Höhe erleichtert die Arbeit, sofern man Polhöhe und Uhrstand schon ungefähr kennt. Kann man von einem Stativ aus beobachten,
\SourceLeaf{136}%
so ist das nicht nur wesentlich bequemer. Man kann auch die Berührung der Bilder viel genauer und besser in der Mitte des Gesichtsfelds erkennen als bei einer Beobachtung aus freier Hand. Mit fleißiger Übung lässt sich aber auch darin viel erreichen.

Ein Nachteil des zuvor beschriebenen Verfahrens ist, dass man nachts bei künstlichem Licht die feinen Teilstriche des Sextanten weniger gut ablesen kann als am Tag. Wichtiger ist jedoch ein Umstand, der erst seit einiger Zeit zur Sprache gekommen ist. Bei vielen Spiegelsextanten, selbst bei solchen von den besten Meistern, ist die Teilung jedenfalls etwas weniger vollkommen, als man ihnen bei uns bisher ohne Weiteres zugetraut hat. Ich selbst besitze einen 10-Zoll-Spiegelsextanten von \textsc{Troughton}, Nr.\,420. Er ist sonst ausgezeichnet, gibt die Winkel von $100$ bis $120^0$ aber ganz entschieden um $50$ bis $60$ oder $70''$ zu klein an. Da alle übrigen Prüfungen und Berichtigungen an ihm mit größter Sorgfalt vorgenommen worden sind, kann ich das nur einem Teilungsfehler zuschreiben. Auf dem Meer sind solche Fehler bei der Breiten- und Zeitbestimmung ohne jede Bedeutung, und an Land fallen sie bei Orten, deren Lage bisher völlig unbekannt war, wenig ins Gewicht. Wo man jedoch größere Genauigkeit erreichen will, darf man die Vollkommenheit der Teilung nicht einfach voraussetzen. Man muss entweder zuerst versuchen, die Teilungsfehler möglichst genau zu ermitteln, oder darauf verzichten, mit Verfahren, die genau gemessene Höhen voraussetzen, völlig zuverlässige Bestimmungen vorzunehmen.

Aus diesen Gründen dürfte vielen Beobachtern ein Verfahren willkommen sein, mit dem man Polhöhe und Uhrstand in kurzer Zeit sehr genau bestimmen kann, selbst wenn das Instrument noch so ungenau oder unter Umständen gar nicht geteilt ist. Auf diese Weise kann man auch etwaige Teilungsfehler des Instruments selbst bestimmen. Die Genauigkeit der Ergebnisse hängt also allein davon ab, wie sorgfältig man die Berührung der Bilder beobachtet. Beim heutigen Stand unserer Sternverzeichnisse dürfen die verwendeten Sternpositionen als vollkommen fehlerfrei gelten. Bei diesem Verfahren wartet man die Zeitpunkte ab, zu denen drei beliebige Sterne in Vertikalkreisen, die am Zenit keine zu spitzen Winkel bilden, dieselbe Höhe erreichen. Diese Höhe ist im Übrigen frei wählbar und braucht
\SourceLeaf{137}%
selbst nicht bekannt zu sein. Polhöhe und Uhrstand lassen sich daraus wohl am bequemsten auf die folgende Weise ableiten. Ich verwende die folgenden Bezeichnungen.
\begin{description}
\item[$\alpha,\alpha',\alpha''$] die Rektaszensionen der drei Sterne,
\item[$\delta,\delta',\delta''$] ihre Deklinationen, wobei südliche Deklinationen als negativ gelten,
\item[$\theta,\theta',\theta''$] die drei Uhrzeiten, zu denen diese Sterne die Höhe $h$ erreichen,
\item[$k$] die Voreilung der Uhr gegenüber der Sternzeit, die ich für alle drei Beobachtungen als gleich annehme. Ginge die Uhr nicht nach Sternzeit, könnte man unter $k$ ihre Voreilung zu einem beliebigen Zeitpunkt verstehen. Mit $\theta,\theta',\theta''$ wären dann die Zeiten gemeint, die die Uhr angezeigt hätte, wenn sie zwischen diesem Zeitpunkt und den Beobachtungen genau nach Sternzeit gegangen wäre. Man könnte etwa $k$ als die Voreilung bei der ersten Beobachtung ansetzen und die beiden anderen Zeiten entsprechend vergrößern oder verkleinern, wenn die Uhr langsamer oder schneller als nach Sternzeit liefe.
\item[$\varphi$] die Polhöhe des Ortes.
\end{description}
Die drei Stundenwinkel sind also offenbar $\theta-k-\alpha$, $\theta'-k-\alpha'$, $\theta''-k-\alpha''$, in Winkel umgerechnet. Setzt man
\[
\begin{aligned}
\theta-\alpha&=t\\
\theta'-\alpha'&=t'\\
\theta''-\alpha''&=t''
\end{aligned}
\]
so erhält man die folgenden drei Gleichungen.
\[
\begin{aligned}
\mathrm{I.}\qquad&\sin h=\sin\varphi\sin\delta+\cos\varphi\cos\delta\cos(t-k)\\
\mathrm{II.}\qquad&\sin h=\sin\varphi\sin\delta'+\cos\varphi\cos\delta'\cos(t'-k)\\
\mathrm{III.}\qquad&\sin h=\sin\varphi\sin\delta''+\cos\varphi\cos\delta''\cos(t''-k)
\end{aligned}
\]

Zieht man I von II ab, ergibt sich nach einer einfachen Umformung
\[
\begin{aligned}
&2\sin\varphi\sin\tfrac12(\delta'-\delta)\cos\tfrac12(\delta'+\delta)\\
&\quad=2\cos\varphi\cos[\tfrac12(t+t')-k]\cos\tfrac12(t'-t)\sin\tfrac12(\delta'-\delta)\sin\tfrac12(\delta'+\delta)\\
&\qquad+2\cos\varphi\sin[\tfrac12(t+t')-k]\sin\tfrac12(t'-t)\cos\tfrac12(\delta'-\delta)\cos\tfrac12(\delta'+\delta)
\end{aligned}
\]
oder
\[
\begin{aligned}
\operatorname{tang}\varphi&=\cos[\tfrac12(t+t')-k]\cos\tfrac12(t'-t)\operatorname{tang}\tfrac12(\delta'+\delta)\\
&\quad+\sin[\tfrac12(t+t')-k]\sin\tfrac12(t'-t)\operatorname{cotg}\tfrac12(\delta'-\delta)
\end{aligned}
\]
\SourceLeaf{138}%
Man bestimme $A'$ und $B'$ so, dass
\[
\begin{aligned}
A'\sin B'&=\sin\tfrac12(t'-t)\operatorname{cotg}\tfrac12(\delta'-\delta)\\
A'\cos B'&=\cos\tfrac12(t'-t)\operatorname{tang}\tfrac12(\delta'+\delta)
\end{aligned}
\]
gilt. Setzt man zugleich
\[
\tfrac12(t+t')-B'=C'
\]
so nimmt die vorige Gleichung die folgende Form an
\[
\mathrm{IV.}\qquad\operatorname{tang}\varphi=A'\cos(C'-k)
\]

Auf genau dieselbe Weise oder einfach, indem man die Größen der zweiten Beobachtung mit denen der dritten vertauscht, erkennt man den entsprechenden Zusammenhang. Bestimmt man $A''$ und $B''$ so, dass
\[
\begin{aligned}
A''\sin B''&=\sin\tfrac12(t''-t)\operatorname{cotg}\tfrac12(\delta''-\delta)\\
A''\cos B''&=\cos\tfrac12(t''-t)\operatorname{tang}\tfrac12(\delta''+\delta)
\end{aligned}
\]
gilt, und setzt man zugleich
\[
\tfrac12(t+t'')-B''=C''
\]
so ergibt sich die folgende Gleichung
\[
\mathrm{V.}\qquad\operatorname{tang}\varphi=A''\cos(C''-k)
\]

Aus IV und V zusammen lassen sich nun leicht $k$ und $\varphi$ ableiten. Es gilt nämlich
\[
A'\cos(C'-k)=A''\cos(C''-k)
\]
also
\[
\begin{aligned}
&(A''-A')\cos[\tfrac12(C'+C'')-k]\cos\tfrac12(C''-C')\\
&=(A'+A'')\sin[\tfrac12(C'+C'')-k]\sin\tfrac12(C''-C')
\end{aligned}
\]
Setzt man also
\[
\frac{A'}{A''}=\operatorname{tang}\zeta
\]
woraus
\[
\frac{A''-A'}{A''+A'}=\operatorname{tang}(45^0-\zeta)
\]
\SourceLeaf{139}%
folgt, und bestimmt $\psi$ aus der Gleichung
\[
\operatorname{tang}(45^0-\zeta)\operatorname{cotang}\tfrac12(C''-C')=\operatorname{tang}\psi
\]
so erhält man
\[
k=\tfrac12(C'+C'')-\psi
\]
und anschließend $\varphi$ aus einer der beiden Gleichungen IV oder V. Aus der so gefundenen Polhöhe und den Stundenwinkeln $t-k$, $t'-k$, $t''-k$ kann man anschließend, wenn man möchte, mit den Formeln I, II, III oder mit anderen bekannten Verfahren die Höhe $h$ ableiten. Alle drei Beobachtungen müssen denselben Wert für sie ergeben.

Zu dieser Lösung sind noch die folgenden Bemerkungen zu machen.

1) Aus den beiden Gleichungen $A\sin B=M$, $A\cos B=N$ bestimmt man die Größen $A$ und $B$ mit den Formeln
\[
\begin{gathered}
\operatorname{tang}B=\frac{M}{N}\\
A=\frac{M}{\sin B}\quad\text{oder}\quad A=\frac{N}{\cos B}
\end{gathered}
\]
Für $A$ nimmt man vorzugsweise den ersten Ausdruck, wenn $M$ größer als $N$ ist, und im entgegengesetzten Fall den anderen. Bestimmt man den Winkel $B$ aus seiner Tangente, so bleibt bei positiver Tangente die Wahl zwischen dem ersten und dritten Quadranten offen, bei negativer zwischen dem zweiten und vierten. Man wählt den Quadranten, in dem der Sinus das Vorzeichen von $M$ und der Kosinus das Vorzeichen von $N$ hat. Dadurch wird $A$ positiv.

2) Die ähnliche Zweideutigkeit bei der Bestimmung von $\psi$ aus der Tangente muss man so auflösen, dass $\operatorname{tang}\varphi$ positiv wird. Man nimmt also $\psi$ zwischen $-90^0$ und $+90^0$, vorausgesetzt, die Beobachtungen wurden auf der Nordhalbkugel gemacht. Auf der Südhalbkugel wäre es umgekehrt.

Bei jedem Verfahren, das man dem praktischen Astronomen zur Anwendung empfiehlt, ist es eine unerlässliche Pflicht, den Einfluss der unvermeidlichen Beobachtungsfehler auf die Ergebnisse abzuschätzen. So kann man sich davon überzeugen, ob und unter welchen Umständen das Verfahren überhaupt zuverlässig anwendbar ist. Auf die Vernachlässigung dieser Pflicht gehen die vielen unausgereiften Einfälle zurück, über deren Nutzlosigkeit die praktischen Astronomen klagen. Eine einfache Untersuchung, die hier der Kürze wegen entfällt,
\SourceLeaf{140}%
zeigt Folgendes. Es seien $\lambda,\lambda',\lambda''$ die Azimute zu den drei Beobachtungen, vom Südpunkt aus in Richtung der täglichen Bewegung gezählt. Ein Fehler von $\Delta$ Zeitsekunden bei der ersten Berührung zieht dann einen Fehler von
\[
\frac{\Delta\sin\lambda\sin\tfrac12(\lambda''+\lambda')}{2\sin\tfrac12(\lambda'-\lambda)\sin\tfrac12(\lambda''-\lambda)}\quad\text{Zeitsekunden}
\]
bei der Zeitbestimmung und von
\[
\frac{15\Delta\sin\lambda\cos\varphi\cos\tfrac12(\lambda''+\lambda')}{2\sin\tfrac12(\lambda'-\lambda)\sin\tfrac12(\lambda''-\lambda)}\quad\text{Bogensekunden}
\]
bei der Polhöhe nach sich. Hat man die Berührung zu früh beobachtet, muss man nämlich $k$ um die eine Größe vermindern und $\varphi$ um die andere erhöhen. Die Ausdrücke für den Einfluss der Fehler bei der zweiten und dritten Beobachtung sind ganz ähnlich und ergeben sich allein durch Vertauschen von $\lambda$ mit $\lambda'$ oder $\lambda''$. Da die Höhe des ersten Sterns in $\Delta$ Zeitsekunden um $15\Delta\sin\lambda\cos\varphi$ Bogensekunden abnimmt, lässt sich der Einfluss der Beobachtungsfehler noch einfacher ausdrücken. Angenommen, die Höhe des ersten Sterns war in dem Augenblick, in dem man sie mit der Einstellung des Instruments übereinstimmend fand, tatsächlich noch oder schon um $D$ Bogensekunden größer. Dann erhält man für $k$ in Bogensekunden einen um
\[
\frac{D\sin\tfrac12(\lambda''+\lambda')}{2\cos\varphi\sin\tfrac12(\lambda'-\lambda)\sin\tfrac12(\lambda''-\lambda)}
\]
zu großen und für $\varphi$ einen um
\[
\frac{D\cos\tfrac12(\lambda''+\lambda)}{2\sin\tfrac12(\lambda'-\lambda)\sin\tfrac12(\lambda''-\lambda)}
\]
zu kleinen Wert.

Daraus folgt, dass man nur darauf achten muss, dass keiner der Sinuswerte von $\tfrac12(\lambda'-\lambda)$, $\tfrac12(\lambda''-\lambda)$, $\tfrac12(\lambda''-\lambda')$ zu klein wird. Dazu wählt man nur Sterne, die die Höhe $h$ bei recht unterschiedlichen Azimuten erreichen. Zweitens ist klar, dass Sterne, deren Höhe sich langsam ändert, ebenso brauchbar sind wie solche, die rasch steigen oder fallen. Bei den ersteren kommt es keineswegs darauf an, den Augenblick haargenau zu treffen, in dem sie die verlangte Höhe erreichen. Entscheidend ist nur, dass sie im dafür angesetzten Augenblick tatsächlich nicht merklich davon abweichen. Man kann daher ohne Bedenken auch Sterne nahe der Kulmination oder den Polarstern wählen. Gerade sie sind sehr geeignet,
\SourceLeaf{141}%
weil man die eben genannte Anforderung in Ruhe erfüllen kann. Wenigstens einer der drei Sterne wird seine Höhe ohnehin immer schneller ändern, wenn die Bedingung der unterschiedlichen Azimute erfüllt ist.

Bevor man diese Beobachtung vornimmt, sollte man sich darauf vorbereiten. Zu jeder Stunde lassen sich leicht einige gut erkennbare Sterne finden, die kurz nacheinander dieselbe Höhe erreichen. Mit einem Globus oder einer stereografischen Projektion der Himmelskugel lässt sich das ohne Weiteres bewerkstelligen. Kennt man die Zeitabstände zwischen den drei Beobachtungen im Voraus auf ein paar Minuten genau, so gelingt die Beobachtung umso besser. Man müsste schlecht beobachtet haben, wenn man damit den Uhrstand nicht wenigstens auf ein paar Sekunden genau erhielte. Durch die Kombination mehrerer Sterne und die Wiederholung an mehreren Abenden, an denen man auch dieselben Höhen erneut messen kann, lässt sich mit diesem Verfahren ein sehr hoher Grad an Genauigkeit erreichen. Außerdem lassen sich mehrere Abkürzungen der Rechnung angeben, besonders wenn man dieselben Sterne öfter in derselben oder nahezu derselben Höhe beobachtet hat. Darauf gehe ich hier nicht näher ein.

Um dieses Verfahren weiter zu erläutern, füge ich hier die vollständige Rechnung für eine tatsächliche Anwendung bei. Am 27. August dieses Jahres stellte ich meinen Sextanten auf die doppelte Höhe $105^0\ 18'\ 55''$ ein und beobachtete mit dem Quecksilberhorizont aus freier Hand die Sterne $\alpha$ Andromeda, $\alpha$ Kleiner Bär und $\alpha$ Leier. An der genau nach Sternzeit gehenden Pendeluhr von \textsc{Shelton} las ich die folgenden Zeiten ab.
\begin{center}
\begin{tabular}{lrrr}
$\alpha$ Andromeda $\ .\ .\ .\ .$&$21^{\mathrm u}$&$33^{\mathrm m}$&$26^{\mathrm s}$\\
$\alpha$ Kleiner Bär $\ .\ .\ .\ .$&$21$&$47$&$30$\\
$\alpha$ Leier $\ .\ .\ .\ .\ .\ .$&$22$&$5$&$21$
\end{tabular}
\end{center}
Die scheinbaren Positionen der Sterne an diesen Tagen sind unten angegeben.
\begin{center}
\small
\begin{tabular}{lcc}
&Rektaszension&Nördl. Deklination\\
$\alpha$ Andromeda&$23^{\mathrm u}\ 58^{\mathrm m}\ 33^{\mathrm s}33$&$28^0\ 2'\ 14''8$\\
$\alpha$ Kleiner Bär&$0\quad55\quad4.7$&$88\quad17\quad5.7$\\
$\alpha$ Leier $\ .\ .$&$18\quad30\quad28.96$&$38\quad37\quad6.6$
\end{tabular}
\end{center}
Daraus ergibt sich Folgendes.
\SourceLeaf{142}%
\begin{center}\small
\begin{tabular}{rcrcl}
&In Zeit&&&Im Winkelmaß\\
$t=$&$21^{\mathrm u}\ 34^{\mathrm m}\ 52^{\mathrm s}67$&$\ .\ .\ .$&&$323^0\ 43'\ 10''05$\\
$t'=$&$20\quad52\quad25.30$&$\ .\ .\ .$&&$313\quad6\quad19.50$\\
$t''=$&$3\quad34\quad52.04$&$\ .\ .\ .$&&$53\quad43\quad0.60$
\end{tabular}
\par\medskip
\begin{minipage}[t]{.49\textwidth}\centering
\begin{tabular}[t]{@{}lr@{}}
$\tfrac12(t'-t)=$&$-5^0\ 18'\ 25''27$\\
$\tfrac12(t'+t)=$&$318\quad24\quad44.77$\\[6pt]
$\tfrac12(\delta'-\delta)=$&$30\quad7\quad25.45$\\
$\tfrac12(\delta'+\delta)=$&$58\quad9\quad40.25$\\[6pt]
$\log\sin\tfrac12(t'-t)\;.\;.$&$8.9661069\,\mathrm n$\\
$\log\cot\tfrac12(\delta'-\delta)\;.\;.$&$0.2363974$\\[6pt]
$\log\cos\tfrac12(t'-t)\;.\;.$&$9.9981343$\\
$\log\operatorname{tang}\tfrac12(\delta'+\delta)\;.\;.$&$0.2069331$
\end{tabular}
\end{minipage}\hfill
\begin{minipage}[t]{.49\textwidth}\centering
\begin{tabular}[t]{@{}lr@{}}
$\tfrac12(t''-t)=$&$-135^0\ 0'\ 4''72$\\
$\tfrac12(t''+t)=$&$188\quad43\quad5.32$\\[6pt]
$\tfrac12(\delta''-\delta)=$&$5\quad17\quad25.90$\\
$\tfrac12(\delta''+\delta)=$&$33\quad19\quad40.70$\\[6pt]
$\log\sin\tfrac12(t''-t)\;.\;.$&$9.8494751\,\mathrm n$\\
$\log\cot\tfrac12(\delta''-\delta)\;.\;.$&$1.0333869$\\[6pt]
$\log\cos\tfrac12(t''-t)\;.\;.$&$9.8494949\,\mathrm n$\\
$\log\operatorname{tang}\tfrac12(\delta''+\delta)\;.\;.$&$9.8179461$
\end{tabular}
\end{minipage}
\end{center}
Damit haben wir
\begin{center}\small
\begin{minipage}[t]{.49\textwidth}\centering
\begin{tabular}[t]{@{}lr@{}}
$\log A'\cos B'\;.\;.\;.$&$0.2050674$\\
$\log A'\sin B'\;.\;.\;.$&$9.2025043\,\mathrm n$
\end{tabular}
\end{minipage}\hfill
\begin{minipage}[t]{.49\textwidth}\centering
\begin{tabular}[t]{@{}lr@{}}
$\log A''\cos B''\;.\;.\;.$&$9.6674410\,\mathrm n$\\
$\log A''\sin B'\;.\;.\;.$&$0.8828620\,\mathrm n$
\end{tabular}
\end{minipage}
\end{center}
Daraus erhalten wir
\begin{center}\small
\begin{minipage}[t]{.49\textwidth}\centering
\begin{tabular}[t]{@{}lr@{}}
$B'=$&$354^0\ 19'\ 22''04$\\
$B''=$&$266\quad30\quad55.07$\\[6pt]
$C'=$&$-35\quad54\quad37.27$\\
$C''=$&$-77\quad47\quad49.75$\\[6pt]
$\tfrac12(C''-C')=$&$-20^0\ 56'\ 36''24$\\
$\tfrac12(C''+C')=$&$-56\quad51\quad13.51$
\end{tabular}
\end{minipage}\hfill
\begin{minipage}[t]{.49\textwidth}\centering
\begin{tabular}[t]{@{}lr@{}}
$\log A'\;.\;.\;.$&$0.2072029$\\
$\log A''\;.\;.\;.$&$0.8836657$\\[6pt]
$\log\operatorname{tang}\zeta\;.\;.$&$9.3235372$\\
$\zeta=$&$11^0\ 53'\ 41''28$\\[6pt]
$45^0-\zeta=$&$33\quad6\quad18.72$
\end{tabular}
\end{minipage}
\par\medskip
\begin{tabular}{lr}
$\log\operatorname{tang}(45^0-\zeta)\;.\;.$&$9.8142617$\\
$\log\cot\tfrac12(C''-C')\;.\;.$&$0.4171063\,\mathrm n$\\\hline
$\log\operatorname{tang}\psi\;.\;.\;.\;.$&$0.2313680\,\mathrm n$\\
\multicolumn{2}{l}{$\psi=-59^0\ 35'\ 14''71$}\\
\multicolumn{2}{l}{$k=+2\quad44\quad1.20$}
\end{tabular}
\end{center}
Die Uhr ging der Sternzeit also um $10^{\mathrm m}\,56^{\mathrm s}08$ voraus.
\SourceLeaf{143}%
Das stimmt auf ein paar Zehntel mit dem Wert überein, der aus den Sterndurchgängen am Mauerquadranten abgeleitet worden war.

Zur Bestimmung der Polhöhe haben wir
\[
\begin{aligned}
C'-k&=-38^0\ 38'\ 38''47\\
C''-k&=-80\quad31\quad50.95
\end{aligned}
\]
\begin{center}\small
\begin{minipage}[t]{.49\textwidth}\centering
\begin{tabular}[t]{@{}lr@{}}
$\log A'\;.\;.\;.\;.$&$0.2072029$\\
$\log\cos(C'-k)\;.$&$9.8926738$\\\hline
$\log\operatorname{tang}\varphi\;.\;.\;.$&$0.0998767$\\
\multicolumn{2}{l}{$\varphi=51^0\ 31'\ 51''51$}
\end{tabular}
\end{minipage}\hfill
\begin{minipage}[t]{.49\textwidth}\centering
\begin{tabular}[t]{@{}lr@{}}
$\log A''\;.\;.\;.\;.$&$0.8836657$\\
$\log\cos(C''-k)\;.$&$9.2162109$\\\hline
&$0.0998766$
\end{tabular}
\end{minipage}
\end{center}

Berechnet man mit diesen Ergebnissen die wahre Höhe der Sterne, so erhält man aus allen drei Sternen bis auf eine Hundertstelsekunde denselben Wert.
\[
h=52^0\ 37'\ 21''3
\]
Die Refraktion beträgt nach der Korrektur anhand des Barometer- und Thermometerstands $42''7$. Die scheinbare Höhe ist also $52^0\ 38'\ 4''$. Der Kollimationsfehler des Sextanten beträgt $-3'\ 30''$, der gemessene Winkel also $105^0\ 15'\ 25''$. Da der wahre Winkel $105^0\ 16'\ 8''$ beträgt, ist der gemessene um $43''$ zu klein.

Am 25. August war der Sextant auf einen um $5''$ größeren Winkel eingestellt, nämlich $105^0\ 19'\ 0''$. Dieselben drei Sterne wurden zu den unten angegebenen Uhrzeiten beobachtet.
\begin{center}
\begin{tabular}{lrrr}
$\alpha$ Andromeda $\ .\ .\ .\ .$&$21^{\mathrm u}$&$33^{\mathrm m}$&$29^{\mathrm s}$\\
$\alpha$ Kleiner Bär $\ .\ .\ .\ .$&$21$&$47$&$38$\\
$\alpha$ Leier $\ .\ .\ .\ .\ .\ .$&$22$&$5$&$22$
\end{tabular}
\end{center}
Daraus ergeben sich die folgenden Werte.
\begin{center}
\begin{tabular}{lr}
Voreilung der Uhr gegenüber der Sternzeit&$10^{\mathrm m}\,57^{\mathrm s}9$\\
$\varphi=$&$51^0\ 31'\ 56''7$\\
$h=$&$52\quad37\quad29.15$
\end{tabular}
\end{center}
Die doppelte scheinbare Höhe beträgt $105^0\ 16'\ 25''5$. Der Sextant zeigt $105^0\ 15'\ 30''$ an, sein Fehler beträgt also $-54''5$.

Im Mittel beträgt der Fehler des Sextanten beim Winkel $105^0$ also $-48''7$, die Polhöhe $51^0\ 31'\ 54''1$. Einige Beobachtungen von Professor \textsc{Harding}
\SourceLeaf{144}%
an denselben Sternen, aber in etwas größerer Höhe, sowie eine frühere Beobachtung von mir mit $\alpha$ Andromeda, $\alpha$ Adler und $\alpha$ Leier ergeben für die Polhöhe bis auf ein paar Sekunden denselben Wert.

Um den Einfluss der Beobachtungsfehler besser überblicken zu können, habe ich aus den Zahlen des hier ausgeführten Beispiels die Azimute zu den drei Sternen berechnet. Ich erhielt die unten angegebenen Werte.
\begin{center}
\begin{tabular}{lcl}
$\alpha$ Andromeda & $.\ .\ .\ .$ & $293^{0}\ 45'\ 15''$\\
$\alpha$ Kleiner Bär & $.\ .\ .\ .$ & $182\quad9\quad9$\\
$\alpha$ Leier & $.\ .\ .\ .\ .\ .$ & $90\quad17\quad52$
\end{tabular}
\end{center}
Daraus finde ich mit den obigen Formeln, dass bei um $\Delta,\Delta',\Delta''$ Zeitsekunden zu früh beobachteten Berührungen die Voreilung der Uhr um
\[
+0.391\Delta+0.0066\Delta'+0.603\Delta''
\]
Zeitsekunden zu klein und die Polhöhe um
\[
+3.808\Delta-0.2884\Delta'-3.519\Delta''
\]
Bogensekunden zu klein ausfällt.

Irrt man sich also bei $\alpha$ Kleiner Bär um $20''$ und bei jedem der beiden anderen Sterne um $1''$, so kann der Fehler der Zeitbestimmung höchstens $1^{\mathrm{s}}1$ und der Fehler der Polhöhe höchstens $13''$ betragen, sofern die Sternpositionen selbst genau und zuverlässig sind. Will man die möglichen Fehler der Endergebnisse aus den Fehlern der Höhen statt aus denen der Zeiten bestimmen, dienen dazu die anderen oben angegebenen Formeln. Nehmen wir an, dass die Sternhöhen in den drei Beobachtungszeitpunkten um $D,D',D''$ Bogensekunden größer waren, als es die Einstellung des Sextanten erforderte. Daraus ergibt sich ein Fehler von
\[
+0.0430D+0.0177D'-0.0607D''
\]
Zeitsekunden bei der Voreilung der Uhr und von
\[
+0.446D-0{,}823D'+0.377D''
\]
Bogensekunden bei der Polhöhe. Beide Werte fallen um den jeweiligen Betrag zu groß aus. Die letzten Formeln dienen zugleich dazu, den Einfluss etwaiger Unsicherheiten der Sternpositionen abzuschätzen.
\SourceLeaf{145}%
Man braucht unter $D,D',D''$ nur die Fehler dieser Positionen in Richtung der drei Vertikalkreise und in Sekunden eines Großkreises zu verstehen, um die gefundenen Formeln unmittelbar auf ihren Einfluss anzuwenden. Denn $D,D',D''$ sind die Unterschiede zwischen den Höhen, die der Einstellung des Sextanten entsprechen, und den Höhen derjenigen Punkte der Himmelskugel, deren Rektaszensionen und Deklinationen der Rechnung zugrunde liegen. Es ist einleuchtend, dass es gleichgültig ist, ob sie bei völliger Übereinstimmung der Sterne mit diesen Punkten aus Beobachtungsfehlern stammen oder bei völlig genauen Beobachtungen aus kleinen Unterschieden zwischen den Sternen und jenen Punkten. Schließlich bemerke ich noch, dass die Summe der Koeffizienten in jeder der beiden Formeln für den Fehler der Polhöhe ebenso wie in der zweiten Formel für den Fehler der Zeitbestimmung stets $=0$ ist. In der ersten Formel für den Fehler der Zeitbestimmung ist sie dagegen $=1$. Den Grund dafür wird man bei einigem Nachdenken leicht finden.

\par\bigskip\centerline{\rule{.27\linewidth}{.4pt}}
